\sectionbrief
{one variable $x$, an ODE $p(D)u=0$ with constant coefficients, $D=d/dx$, $p$ known.}
{in one variable every object of the later sections (minimal annihilator, spectral projectors, companion matrix) appears with explicit formulas and no geometry of curves.}
{prototype for Section~\ref{sec:ident} (identification) and Section~\ref{sec:groups} (the Lagrange projectors reappear as Eq.~\eqref{xeq:lagrange}).}
This section introduces, in the simplest setting, objects that reappear in Sections~\ref{sec:ident} and~\ref{sec:groups}: the minimal annihilator, the companion matrix and the spectral projectors.

\begin{notebox}{Setting and symbols}
\begin{itemize}[nosep]
\item We consider an ODE $p(D)u=0$ in one variable $x$, with $D=d/dx$ (in this section only; elsewhere $D$ is a diffusion constant).
\item The polynomial $p(\lambda)=\prod_j(\lambda-\lambda_j)^{m_j}$ has distinct roots $\lambda_j\in\C$ with multiplicities $m_j\ge1$, so that $m=\deg p=\sum_jm_j$ is the \emph{order} of the equation (not to be confused with the Sobolev index $m$ of Section~\ref{sec:groups}).
\item The set of solutions is the \emph{kernel} $\ker p(D)$, and
\[
\ker p(D)=\Big\{\sum_jq_j(x)e^{\lambda_jx}:\ \deg q_j<m_j\Big\},
\]
where each $q_j$ is a polynomial of degree $<m_j$.
\end{itemize}
\end{notebox}

\begin{aside}{Why this form}
For one root, $(D-\lambda)\big(q\,e^{\lambda x}\big)=q'e^{\lambda x}$: applying $D-\lambda$ lowers the degree of $q$ by one. Hence $(D-\lambda_j)^{m_j}$ kills $q\,e^{\lambda_jx}$ exactly when $q^{(m_j)}=0$, i.e.\ $\deg q<m_j$.

Each root therefore contributes the $m_j$ functions $x^ke^{\lambda_jx}$ ($0\le k<m_j$): these are the natural basis functions of the equation, finitely many, and together they form a basis of $\ker p(D)$, a space of dimension $m$.
\end{aside}

\begin{aside}{Example}
$p(\lambda)=(\lambda-1)^2(\lambda+2)^3$ has $\lambda_1=1$, $m_1=2$, $\lambda_2=-2$, $m_2=3$, order $m=5$. A basis of the solutions is
\[
e^{x},\quad xe^{x},\quad e^{-2x},\quad xe^{-2x},\quad x^2e^{-2x},
\]
so every solution is
\[
c_1e^{x}+c_2xe^{x}+c_3e^{-2x}+c_4xe^{-2x}+c_5x^2e^{-2x}.
\]
\end{aside}

\subsection{Identification by the Wronskian}
\begin{proposition}[{(ODE identification)}]
Let $u\in\ker p(D)$. The minimal annihilator of $u$ has degree equal to the number of independent basis functions actually present in $u$; it equals $p$ if and only if $u,u',\dots,u^{(m-1)}$ are linearly independent (nonvanishing Wronskian), i.e.\ iff every root of $p$ appears in $u$ with its full multiplicity.
\end{proposition}
\emph{Reason.} $\dimc\spn\{u,u',u'',\dots\}$ is the degree of the minimal annihilator (Lemma~\ref{lem:ep} in one variable); a missing root lowers it.

\subsection{The companion matrix and spectral form}
Write $p(\lambda)=\lambda^m+a_{m-1}\lambda^{m-1}+\dots+a_0$ (monic) and put $X=(u,u',\dots,u^{(m-1)})^{\top}$. Then $X'=CX$ with
\[
C=\begin{pmatrix}0&1&&\\&0&\ddots&\\&&\ddots&1\\-a_0&-a_1&\cdots&-a_{m-1}\end{pmatrix},\qquad
X(x)=e^{xC}X(0).
\]
The eigenvalues of $C$ are the roots of $p$. For distinct roots the matrix exponential has the spectral form
\begin{equation}\label{eq:spectral}
\begin{aligned}
e^{xC}&=\sum_{j=1}^me^{\lambda_jx}\,\Pi_j,\\
\Pi_j&=\prod_{k\ne j}\frac{C-\lambda_k\Id}{\lambda_j-\lambda_k},
\end{aligned}
\end{equation}
where the $\Pi_j$ are the Lagrange projectors onto the eigenspaces. The last row of $C$ carries the whole equation; $C$ is a \emph{matrix over the base field} for an ODE, but, as soon as the coefficients depend on $x$, or there is a second variable, its last row has entries that are operators and the object becomes a matrix over a noncommutative or polynomial ring. The invariant statement is that the whole equation corresponds to the cyclic module $\ring/\Ann(u)$ (Section~\ref{sec:theory}); the companion matrix is a way to compute with it relative to one chosen direction.

\begin{proposition}[{(Projectors of the companion matrix) \led{II-T2} \proved[]}]
The projectors $\Pi_j$ are explicit rational functions of the coefficients, with poles exactly where roots collide; they are polynomial in the coefficients if and only if the corresponding varieties are disjoint. They are \emph{not} orthogonal in general: the companion matrix is not normal for $m=2$ \falsest{} (the contrary was an early conjecture).
\end{proposition}
Hence Hilbert-space orthogonality arguments need an inner product adapted to the problem, not the standard one on the coefficient vector.

\subsection{Tensor construction versus module construction}
One may encode a PDE in a higher-order tensor of derivatives, $J=\partial_{x_{i_1}}\otimes\cdots\otimes\partial_{x_{i_n}}$, or by the companion form. The comparison, \led{Gate 3}: the \emph{principal symbol tensor} $\sigma_m(P)\in\mathrm{Sym}^m(\C^2)$ only sees the highest-order part of $P$ (for the heat operator it sees $\xi^2$ alone), so it is an invariant ``at infinity''; the \emph{cyclic module} $\ring/\Ann(u)$ is the complete invariant. All subsequent results use the module.

\begin{keypoint}
This section shows how an ODE's polynomial encodes its solution space. Its roots determine the exponential rates, their multiplicities determine the polynomial prefactors, the companion matrix represents differentiation, and spectral projectors isolate the different modes. The later sections generalize these ideas to multiple commuting derivatives, where the same structure survives but requires more algebraic machinery: the minimal annihilator becomes the ideal $\Ann(u)$ (Section~\ref{sec:ident}) and the Lagrange projectors reappear as Eq.~\eqref{xeq:lagrange} (Section~\ref{sec:groups}).
\end{keypoint}

%======================================================================
